\documentclass{article}%
\usepackage[T1]{fontenc}%
\usepackage[utf8]{inputenc}%
\usepackage{lmodern}%
\usepackage{textcomp}%
\usepackage{lastpage}%
\usepackage{geometry}%
\geometry{margin=1in}%
%
\usepackage{subcaption}%
\usepackage{amsmath}%
\usepackage{amssymb}%
\usepackage{booktabs}%
\usepackage{longtable}%
\usepackage{graphicx}%
\usepackage{float}%
\usepackage{caption}%
\usepackage{hyperref}%
\usepackage{listings}%
\usepackage{xcolor}%
\lstdefinestyle{mystyle}{backgroundcolor=\color{lightgray}, commentstyle=\color{green}, keywordstyle=\color{blue}, numberstyle=\tiny\color{gray}, stringstyle=\color{red}, basicstyle=\ttfamily\footnotesize, breakatwhitespace=false, breaklines=true, captionpos=b, keepspaces=true, numbers=left, numbersep=5pt, showspaces=false, showstringspaces=false, showtabs=false, tabsize=2}%
\lstset{style=mystyle}%
\title{CHE 565 {-}{-} Exam 2}%
\author{Soki Sem}%
\date{\today}%
%
\begin{document}%
\normalsize%
\setlength{\parindent}{0pt}%
\setlength{\parskip}{0.6em}%
\maketitle%
\section*{Problem 1}%
\label{sec:Problem1}%
\subsection*{Parts A and B: Basic PI Feedback Control}%
\label{subsec:PartsAandBBasicPIFeedbackControl}%
\begin{figure}[H]%
\centering%
\begin{subfigure}{0.45\textwidth}%
\centering%
\includegraphics[width=\linewidth]{Exam2_block_diagram_setpoint.png}%
\subcaption{Setpoint step response}%
\end{subfigure}%
\begin{subfigure}{0.45\textwidth}%
\centering%
\includegraphics[width=\linewidth]{Exam2_block_diagram_disturbance.png}%
\subcaption{Disturbance step response}%
\end{subfigure}%
\end{figure}%
The actual boiler liquid{-}level process is used for the closed{-}loop simulation, while the integrating{-}plus{-}delay model is used only to obtain the controller settings.%
\par%
\begin{align*}%
G_{\text{actual}}(s)=\frac{0.8}{s(2s+1)(5s+1)} \\%
G_{\text{model}}(s)=\frac{0.75}{s}e^{-4s} \\%
G_{\text{d}}(s)=\frac{1}{s+1}e^{-s}%
\end{align*}%
Using the lambda{-}tuning rule with an integrating{-}plus{-}delay model:%
\par%
\begin{align*}%
\lambda = 8\theta \\%
\tau_I = \max(4\theta,\ 2\lambda+\theta) \\%
K' = \frac{K_p}{\tau_p} \\%
K_c = \frac{\tau_I}{K'(\lambda+\theta)^2}%
\end{align*}%
\begin{align*}%
\theta = 4.0 \\%
\lambda = 32.0000 \\%
\tau_I = 68.0000 \\%
K_c = 0.0700%
\end{align*}%
\begin{figure}[H]%
\centering%
\begin{subfigure}{0.45\textwidth}%
\centering%
\includegraphics[width=\linewidth]{problem1_basic_setpoint.png}%
\subcaption{Setpoint step}%
\end{subfigure}%
\begin{subfigure}{0.45\textwidth}%
\centering%
\includegraphics[width=\linewidth]{problem1_basic_disturbance.png}%
\subcaption{Disturbance step}%
\end{subfigure}%
\caption{Basic PI closed{-}loop responses.}%
\end{figure}%
The IAE values for the basic PI controller are:%
\par%
\begin{align*}%
IAE_\text{setpoint,basic} = 31.9746 \\%
IAE_\text{disturbance,basic} = 971.9761%
\end{align*}%
For the setpoint response, the error is measured relative to a final desired value of 1. For the disturbance{-}rejection response, the desired value is 0 because the controller should reject the disturbance and return the output to its original operating point.%
\par

%
\subsection*{Part C: Cascade Control Structure}%
\label{subsec:PartCCascadeControlStructure}%


\begin{figure}[H]%
\includegraphics[width=0.95\textwidth,keepaspectratio]{problem1_cascade_diagram.png}%
\caption{Cascade control diagram.}%
\end{figure}

%
A cascade loop is added by splitting the process into an inner process and an outer process. The inner controller is P{-}only, while the outer controller remains PI. The disturbance is assumed to enter the inner{-}loop process variable, allowing the inner loop to reject the disturbance before it propagates through the slower boiler{-}level dynamics.%
\par%
\begin{align*}%
G_{p1}(s)=\frac{1}{2s+1} \\%
G_{p2}(s)=\frac{0.8}{s(5s+1)} \\%
G_{actual}(s)=G_{p1}(s)G_{p2}(s)%
\end{align*}

%
\subsection*{Part D: Cascade Control with Inner Gain of 50}%
\label{subsec:PartDCascadeControlwithInnerGainof50}%
\begin{figure}[H]%
\centering%
\begin{subfigure}{0.45\textwidth}%
\centering%
\includegraphics[width=\linewidth]{problem1_cascade_setpoint.png}%
\subcaption{Cascade setpoint step}%
\end{subfigure}%
\begin{subfigure}{0.45\textwidth}%
\centering%
\includegraphics[width=\linewidth]{problem1_cascade_disturbance.png}%
\subcaption{Cascade disturbance step}%
\end{subfigure}%
\caption{Cascade closed{-}loop responses with inner controller gain equal to 50.}%
\end{figure}%
The IAE values for the cascade controller are:%
\par%
\begin{align*}%
K_\text{c,inner} = 50.0000 \\%
IAE_\text{setpoint,cascade} = 30.2838 \\%
IAE_\text{disturbance,cascade} = 19.4403%
\end{align*}

%
\subsection*{Part E: Discussion}%
\label{subsec:PartEDiscussion}%
The effectiveness of the cascade controller can be evaluated by comparing the IAE values for the basic feedback controller and the cascade controller. A lower IAE indicates better overall control performance because the total accumulated error is smaller.%
\par%
\begin{align*}%
IAE_\text{setpoint,basic} = 31.9746 \\%
IAE_\text{setpoint,cascade} = 30.2838 \\%
IAE_\text{disturbance,basic} = 971.9761 \\%
IAE_\text{disturbance,cascade} = 19.4403%
\end{align*}%
The cascade controller improves disturbance rejection because the inner loop reacts directly to the disturbance before the disturbance fully propagates through the outer liquid{-}level dynamics.%
\par

%
\section*{Problem 2}%
\label{sec:Problem2}%


\begin{figure}[H]%
\includegraphics[width=0.95\textwidth,keepaspectratio]{problem2_feedforward_diagram.png}%
\caption{Feedforward plus feedback control diagram.}%
\end{figure}

%
For the feedforward design, the disturbance is assumed to enter on the plant{-}input side. The manipulated variable and the disturbance therefore both pass through the same boiler{-}level process dynamics.%
\par%
\begin{align*}%
Y(s)=G_p(s)\left[M(s)+G_d(s)D(s)\right] \\%
M(s)=M_{fb}(s)+M_{ff}(s) \\%
M_{ff}(s)=G_{ff}(s)D(s)%
\end{align*}%
For perfect disturbance cancellation, the disturbance contribution inside the plant input must be zero:%
\par%
\begin{align*}%
G_{ff}(s)D(s)+G_d(s)D(s)=0 \\%
G_{ff}(s)=-G_d(s)%
\end{align*}%
Using the disturbance transfer function:%
\par%
\begin{align*}%
G_d(s)=\frac{1}{s+1}e^{-s} \\%
\boxed{G_{ff}(s)=-\frac{1}{s+1}e^{-s}}%
\end{align*}%
This result follows because the disturbance enters at the same summing point as the manipulated variable before the plant. Therefore, the plant transfer function cancels out in the feedforward design.%
\par

%
\section*{Problem 3}%
\label{sec:Problem3}%
\subsection*{Part A: RGA and Pairing}%
\label{subsec:PartARGAandPairing}%
The steady{-}state gain matrix is found by evaluating each transfer function at s = 0.%
\par%
\begin{align*}%
K=\begin{bmatrix}2.3 & 0.47\\-1.2 & 0.58\end{bmatrix} \\%
\Lambda = K \circ (K^{-1})^T%
\end{align*}%
\begin{align*}%
\Lambda = \begin{bmatrix}0.7028 & 0.2972\\0.2972 & 0.7028\end{bmatrix}%
\end{align*}%
The diagonal RGA values are closer to 1 than the off{-}diagonal values, so the appropriate pairings are:%
\par%
\begin{align*}%
Y_1 \leftrightarrow M_1 \\%
Y_2 \leftrightarrow M_2%
\end{align*}

%
\subsection*{Part B: MIMO Feedback Block Diagram}%
\label{subsec:PartBMIMOFeedbackBlockDiagram}%


\begin{figure}[H]%
\includegraphics[width=0.95\textwidth,keepaspectratio]{problem3_mimo_diagram.png}%
\caption{MIMO feedback control diagram without decouplers.}%
\end{figure}

%
The MIMO process contains four transfer functions. Each manipulated variable affects both output variables, which creates loop interaction.%
\par%
\begin{align*}%
G_{11}(s)=\frac{2.3}{4.6s+1}e^{-0.2s} \\%
G_{12}(s)=\frac{0.47}{2.2s+1}e^{-0.3s} \\%
G_{21}(s)=\frac{-1.2}{18s+1}e^{-0.4s} \\%
G_{22}(s)=\frac{0.58}{1.8s+1}e^{-0.5s}%
\end{align*}

%
\subsection*{Part C: Decoupler Design}%
\label{subsec:PartCDecouplerDesign}%
For diagonal pairing, the decouplers are selected to cancel the off{-}diagonal process interactions.%
\par%
\begin{align*}%
D_{12}(s)=-\frac{G_{12}(s)}{G_{11}(s)} \\%
D_{21}(s)=-\frac{G_{21}(s)}{G_{22}(s)}%
\end{align*}%
Therefore:%
\par%
\begin{align*}%
D_{12}(s)=(-0.2043)\frac{4.6s+1}{2.2s+1}e^{-0.1s} \\%
D_{21}(s)=(2.0690)\frac{1.8s+1}{18s+1}e^{0.1s}%
\end{align*}%
The first decoupler contains an additional delay and is causal. The second decoupler contains a time advance, represented by the positive exponential term. A pure time advance is noncausal, so an exact implementation of this decoupler is not physically realizable. In practice, this decoupler would need to be approximated or modified.%
\par%


\begin{figure}[H]%
\includegraphics[width=0.95\textwidth,keepaspectratio]{problem3_decoupled_mimo_diagram.png}%
\caption{MIMO feedback control diagram with decouplers.}%
\end{figure}

%
\end{document}